% =============================================================================
%  APÊNDICE A: EXERCÍCIOS RESOLVIDOS
% =============================================================================
\chapter{Exercícios Resolvidos}
\label{apend:exercicios}

Este apêndice apresenta soluções detalhadas para exercícios selecionados dos
Capítulos~1~a~8. Cada solução segue o padrão \sdd{}~--~\tdd: definição do
problema, especificação da solução, implementação e validação. Os exercícios
estão organizados por nível de dificuldade, do zero ao PhD.

% =============================================================================
%  SEÇÃO A.1: NÍVEL ZERO
% =============================================================================
\section{Exercícios Nível Zero}
\label{sec:exer-nivel-zero}
\nivelzero

\begin{exercicio}[Lógica Proposicional -- Tabela Verdade]\label{exer:resol-logic}
Construa a tabela verdade da proposição \(P \land (Q \lor \neg R)\).
\end{exercicio}

\textbf{Solução:} A tabela verdade é construída enumerando todas as combinações
de valores lógicos de \(P\), \(Q\) e \(R\) (8 linhas) e computando a expressão
passo a passo:

\[
\begin{array}{ccc|c|c|c}
P & Q & R & \neg R & Q\lor\neg R & P\land(Q\lor\neg R) \\
\midrule
V & V & V & F & V & V \\
V & V & F & V & V & V \\
V & F & V & F & F & F \\
V & F & F & V & V & V \\
F & V & V & F & V & F \\
F & V & F & V & V & F \\
F & F & V & F & F & F \\
F & F & F & V & V & F \\
\end{array}
\]

A proposição é verdadeira quando \(P\) é verdadeiro e ao menos uma das
condições \(Q\) ou \(\neg R\) é verdadeira. Este resultado fundamenta a
construção de guardas lógicas no \texttt{BehavioralGate} do Trust Engine
(Capítulo~5).

\begin{exercicio}[Teoria dos Conjuntos -- Operações]\label{exer:resol-conj}
Dados \(A = \{1,2,3,4\}\), \(B = \{3,4,5,6\}\), calcule \(A \cup B\),
\(A \cap B\) e \(A \setminus B\).
\end{exercicio}

\textbf{Solução:}
\begin{align*}
A \cup B &= \{1,2,3,4,5,6\} \\
A \cap B &= \{3,4\} \\
A \setminus B &= \{1,2\}
\end{align*}

A interseção \(A \cap B = \{3,4\}\) representa os elementos comuns entre
os dois conjuntos, analogamente aos serviços compartilhados entre MCPs
no \opencodeeco{}.

% =============================================================================
%  SEÇÃO A.2: NÍVEL BÁSICO
% =============================================================================
\section{Exercícios Nível Básico}
\label{sec:exer-nivel-basico}
\nivelbasico

\begin{exercicio}[Regressão Linear Simples]\label{exer:resol-regressao}
Dados os pontos \((1,2)\), \((2,3)\), \((3,5)\), \((4,4)\), encontre a reta
de regressão linear \(y = \beta_0 + \beta_1 x\) pelo método dos mínimos
quadrados.
\end{exercicio}

\textbf{Solução:} Calculamos as médias e os somatórios:
\[
\bar{x} = 2{,}5,\quad \bar{y} = 3{,}5,\quad
\sum (x_i - \bar{x})(y_i - \bar{y}) = 4,\quad
\sum (x_i - \bar{x})^2 = 5
\]

\[
\beta_1 = \frac{\sum (x_i - \bar{x})(y_i - \bar{y})}{\sum (x_i - \bar{x})^2}
       = \frac{4}{5} = 0{,}8
\]
\[
\beta_0 = \bar{y} - \beta_1 \bar{x} = 3{,}5 - 0{,}8 \times 2{,}5 = 1{,}5
\]

A reta de regressão é \(\hat{y} = 1{,}5 + 0{,}8x\). Este modelo é
equivalente ao utilizado pelo \texttt{TrustScorer} para predizer a
confiança futura de um agente com base em seu histórico.

% =============================================================================
%  SEÇÃO A.3: NÍVEL INTERMEDIÁRIO
% =============================================================================
\section{Exercícios Nível Intermediário}
\label{sec:exer-nivel-intermediario}
\nivelintermediario

\begin{exercicio}[Autovalores e Autovetores]\label{exer:resol-autoval}
Encontre os autovalores e autovetores da matriz \(A =
\begin{pmatrix} 4 & 1 \\ 2 & 3 \end{pmatrix}\).
\end{exercicio}

\textbf{Solução:} O polinômio característico é:
\[
\det(A - \lambda I) = \det\begin{pmatrix} 4-\lambda & 1 \\ 2 & 3-\lambda \end{pmatrix}
= (4-\lambda)(3-\lambda) - 2 = \lambda^2 - 7\lambda + 10 = 0
\]

Os autovalores são \(\lambda_1 = 5\) e \(\lambda_2 = 2\).

Para \(\lambda_1 = 5\):
\[
(A - 5I)v = \begin{pmatrix} -1 & 1 \\ 2 & -2 \end{pmatrix}v = 0
\Rightarrow v_1 = \begin{pmatrix} 1 \\ 1 \end{pmatrix}
\]

Para \(\lambda_2 = 2\):
\[
(A - 2I)v = \begin{pmatrix} 2 & 1 \\ 2 & 1 \end{pmatrix}v = 0
\Rightarrow v_2 = \begin{pmatrix} 1 \\ -2 \end{pmatrix}
\]

Autovetores são utilizados no \opencodeeco{} para análise de componentes
principais (PCA) na redução de dimensionalidade de embeddings de agentes.

% =============================================================================
%  SEÇÃO A.4: NÍVEL AVANÇADO
% =============================================================================
\section{Exercícios Nível Avançado}
\label{sec:exer-nivel-avancado}
\nivelavancado

\begin{exercicio}[Entropia de Shannon]\label{exer:resol-entropia}
Uma fonte emite símbolos com probabilidades \(p = [0{,}5, 0{,}25, 0{,}15,
0{,}1]\). Calcule a entropia da fonte e o comprimento médio de um código
de Huffman.
\end{exercicio}

\textbf{Solução:} A entropia é:
\[
H(X) = -\sum_{i=1}^{4} p_i \log_2 p_i
\]

\[
H = -(0{,}5 \times \log_2 0{,}5 + 0{,}25 \times \log_2 0{,}25
     + 0{,}15 \times \log_2 0{,}15 + 0{,}1 \times \log_2 0{,}1)
\]

\[
H = -(0{,}5 \times (-1) + 0{,}25 \times (-2) + 0{,}15 \times (-2{,}737)
     + 0{,}1 \times (-3{,}322)) \approx 1{,}743 \text{ bits/símbolo}
\]

Construindo a árvore de Huffman:
\begin{enumerate}
\item Símbolos ordenados: \([0{,}5,\; 0{,}25,\; 0{,}15,\; 0{,}1]\)
\item Combinar \(0{,}15\) e \(0{,}1 \to 0{,}25\)
\item Combinar \(0{,}25\) e \(0{,}25 \to 0{,}5\)
\item Combinar \(0{,}5\) e \(0{,}5 \to 1{,}0\)
\end{enumerate}

Códigos: \(s_1: 0\) (1 bit), \(s_2: 10\) (2 bits), \(s_3: 110\) (3 bits),
\(s_4: 111\) (3 bits). Comprimento médio:
\[
L = 0{,}5 \times 1 + 0{,}25 \times 2 + 0{,}15 \times 3 + 0{,}1 \times 3
  = 1{,}75 \text{ bits/símbolo}
\]

Eficiência: \(\eta = H/L \approx 99{,}6\%\). A compressão de dados no
\emph{Structural Noise Scanner} (SPEC-037) utiliza princípios análogos.

\begin{exercicio}[Gradiente Descendente]\label{exer:resol-grad}
Minimize a função \(f(x) = x^4 - 4x^2 + 4\) usando gradiente descendente
com \(\alpha = 0{,}1\) e \(x_0 = 3\). Realize 3 iterações.
\end{exercicio}

\textbf{Solução:} O gradiente é \(f'(x) = 4x^3 - 8x\).

Iteração 1: \(x_1 = 3 - 0{,}1 \times (4 \times 27 - 24) = 3 - 0{,}1 \times 84 = 2{,}16\)

Iteração 2: \(x_2 = 2{,}16 - 0{,}1 \times (4 \times 10{,}078 - 17{,}28)
          = 2{,}16 - 0{,}1 \times 23{,}032 \approx 1{,}857\)

Iteração 3: \(x_3 = 1{,}857 - 0{,}1 \times (4 \times 6{,}403 - 14{,}856)
          = 1{,}857 - 0{,}1 \times 10{,}756 \approx 1{,}781\)

A função converge para o mínimo local em \(x = \sqrt{2} \approx 1{,}414\).
Este método é empregado pelo \texttt{EvolutionaryTrajectoriesScanner} para
otimizar rotas evolutivas.

% =============================================================================
%  SEÇÃO A.5: NÍVEL PhD
% =============================================================================
\section{Exercícios Nível PhD}
\label{sec:exer-nivel-phd}
\nivelphd

\begin{exercicio}[Teorema de Bayes Aplicado ao Trust Engine]\label{exer:resol-bayes}
Um agente no \opencodeeco{} é classificado como confiável (\(C\)) ou
não-confiável (\(\neg C\)). O \texttt{TrustScorer} emite alertas (\(A\)
quando detecta anomalias). Sabe-se que \(P(C) = 0{,}85\),
\(P(A \mid \neg C) = 0{,}95\) e \(P(A \mid C) = 0{,}05\). Calcule
\(P(C \mid A)\), a probabilidade de um agente ser confiável dado que
um alerta foi emitido.
\end{exercicio}

\textbf{Solução:} Pelo Teorema de Bayes:
\[
P(C \mid A) = \frac{P(A \mid C)P(C)}{P(A \mid C)P(C) + P(A \mid \neg C)P(\neg C)}
\]

Substituindo:
\[
P(C \mid A) = \frac{0{,}05 \times 0{,}85}{0{,}05 \times 0{,}85 + 0{,}95 \times 0{,}15}
= \frac{0{,}0425}{0{,}0425 + 0{,}1425}
= \frac{0{,}0425}{0{,}185} \approx 0{,}2297
\]

Interpretação: quando o \texttt{TrustScorer} emite um alerta, a confiança
no agente cai de 85\% para aproximadamente 23\%. Este cálculo fundamenta
a política de \destaque{shadow mode}: ao detectar anomalia, o agente entra
em modo monitorado até que a confiança seja restabelecida.

\begin{exercicio}[Convergência do Modelo N3 de Auto-Consciência]\label{exer:resol-n3}
Demonstre que o modelo N3 de auto-consciência do \opencodeeco{} é
um ponto fixo da função de auto-observação \(f(M) = M \cup
\texttt{Scan}(M)\), onde \(\texttt{Scan}\) é o operador de varredura
epistemológica.
\end{exercicio}

\textbf{Solução:} Seja \(M_0\) o modelo inicial de si mesmo. Definimos:
\[
M_{t+1} = f(M_t) = M_t \cup \texttt{Scan}(M_t)
\]

Afirmamos que a sequência \((M_t)\) converge para um ponto fixo
\(M^*\) tal que \(f(M^*) = M^*\).

Prova: (1) \(M_t \subseteq M_{t+1}\) por construção (união monótona).
(2) O conjunto de todos os possíveis componentes do ecossistema é
finito (227 skills, 46 MCPs, 128 agentes, etc.). (3) Toda sequência
crescente em um conjunto finito atinge o supremo em tempo finito.
(4) No supremo \(M^*\), \(\texttt{Scan}(M^*) \subseteq M^*\), logo
\(f(M^*) = M^* \cup \texttt{Scan}(M^*) = M^*\).

No \opencodeeco{}, o N3 é atingido quando o \texttt{SelfModel}
(SPEC-036) incorpora todos os componentes detectáveis pelos scanners
epistemológicos. O Behavioral Gate (N3.5) adiciona uma barreira
preventiva sobre este ponto fixo.

\begin{exercicio}[Complexidade do MCSP]\label{exer:resol-mcsp}
O \texttt{MCSP-Solver} (SPEC-032) precisa selecionar o conjunto mínimo
de capacidades para implementar 5 funcionalidades. Cada capacidade
\(c_i\) tem custo \(\text{custo}(c_i)\) e cada funcionalidade \(f_j\)
requer um subconjunto de capacidades. Formule o problema como
otimização inteira.
\end{exercicio}

\textbf{Solução:} Seja \(x_i \in \{0,1\}\) indicando se a capacidade
\(c_i\) é selecionada. O problema é:

\[
\min \sum_{i=1}^{n} \text{custo}(c_i) \cdot x_i
\]

Sujeito a:
\[
\sum_{i: c_i \in R_j} x_i \geq 1,\quad \forall j = 1,\dots,m
\]
\[
x_i \in \{0,1\}, \quad \forall i = 1,\dots,n
\]

Onde \(R_j\) é o conjunto de capacidades requeridas pela funcionalidade
\(f_j\). Este é o problema clássico de \emph{set cover}, NP-difícil.
O \texttt{MCSP-Solver} utiliza um algoritmo guloso com aproximação
\(O(\ln n)\), combinado com poda por \texttt{construction\_cost} real
e desconto por capacidades compartilhadas (SPEC-033).

A solução gulosa ordena capacidades por \(\text{custo}/\text{cobertura}\)
e seleciona iterativamente a de menor custo marginal até cobrir todas
as funcionalidades. Para 5 funcionalidades e 8 capacidades candidatas,
a solução ótima é encontrada por branch-and-bound com poda, tipicamente
em menos de 100 iterações no ecossistema real.
